% !TEX program = xelatex
% SafeBoxForest v4  T-RO 期刊扩展版（中文版）
% 编译: xelatex -> bibtex -> xelatex x 2
\documentclass[journal]{IEEEtran}

% --- 宏包 ---
\usepackage[utf8]{inputenc}
\usepackage{amsmath,amssymb,amsthm}
\usepackage{graphicx}
\usepackage{algorithm}
\usepackage[noend]{algpseudocode}
\def\Or{\textbf{or}}
\newcommand{\alAnd}{\textbf{and}}
\usepackage{booktabs}
\usepackage{multirow}
\usepackage{cite}
\usepackage{xcolor}
\usepackage{hyperref}
\usepackage{cleveref}
\usepackage{bm}
\usepackage{tabularx}
\usepackage{xspace}
\usepackage{microtype}

% --- 中文支持 ---
\usepackage{xeCJK}
\setmainfont{Times New Roman}    % 西文正文，避免 IEEEtran 的 ptm 形状警告
\setsansfont{Arial}
\setmonofont{Courier New}
\setCJKmainfont[BoldFont=SimHei, ItalicFont=KaiTi, BoldItalicFont=KaiTi, SmallCapsFont=SimSun]{SimSun} % 宋体（正文）
\setCJKsansfont{SimHei}          % 黑体（标题）
\setCJKmonofont{FangSong}        % 仿宋（代码）

% --- 定理环境 ---
\newtheorem{theorem}{定理}
\newtheorem{definition}{定义}
\newtheorem{remark}{注}
\newtheorem{proposition}{命题}
\newtheorem{assumption}{假设}

% --- 自定义命令 ---
\newcommand{\Cfree}{\mathcal{C}_\text{free}}
\newcommand{\Cspace}{\mathcal{C}}
\newcommand{\AABB}{\text{AABB}}
\newcommand{\iAABB}{\text{iAABB}}
\newcommand{\FK}{\text{FK}}
\newcommand{\SBF}{\text{SBF}}
\newcommand{\Real}{\mathbb{R}}
\newcommand{\calO}{\mathcal{O}}
\newcommand{\calA}{\mathcal{A}}
\newcommand{\calF}{\mathcal{F}}
\newcommand{\calR}{\mathcal{R}}
\newcommand{\bT}{\mathbf{T}}
\newcommand{\TODO}[1]{\textit{[待完成: #1]}}
\newcommand{\LinkIAABBDef}{由每条连杆的一对端点 iAABB 先取包围盒、再按连杆半径膨胀得到；若启用细分，则表示各细分子段盒的并集。}
\DeclareMathOperator{\interior}{int}

% ==============================================================================
\begin{document}

\title{SafeBoxForest：基于终身包络缓存的\\
快速位形空间盒覆盖与 GCS 运动规划}

\author{
  田宇，任鸿亮
  \thanks{香港中文大学电子工程系，中国香港特别行政区。
  \{ty1997, hlren\}@ee.cuhk.edu.hk}
}

\markboth{IEEE Transactions on Robotics, Vol.~X, No.~Y, 2026}{}

\maketitle

% ==============================================================================
% 摘要
% ==============================================================================
\begin{abstract}
本文提出 SafeBoxForest（SBF）系统，
该系统能够快速将关节机器人的无碰撞位形空间（$\Cfree$）
覆盖为凸盒，
以支持凸集图（Graph-of-Convex-Sets，GCS）运动规划。
SBF 引入了模块化的两层流水线架构：
四种可互换的\emph{端点 iAABB 源}
（IFK、CritSample、AnalyticalCrit、GCPC）
馈入三种\emph{包络表示}
（LinkIAABB\footnotemark、LinkIAABB\_Grid、Hull16\_Grid）。
所有源均产生统一的\emph{端点 iAABB}中间表示——
沿运动学链各端点的区间轴对齐包围盒——
从而将源计算与包络栅格化解耦。
核心组件包括：
(i)~具有 DH 链直接系数提取（零 FK 优化）和
区间正运动学认证的\emph{五阶段解析临界点求解器}
（对面-$k$、$k \leq 2$ 具有完备性）；
(ii)~\emph{几何临界点缓存}（GCPC），
采用八阶段查询流水线，具有 $k\pi/2$ 背景缓存、
增量 DH 链重算和 AA 剪枝，
预制表正运动学评估的临界位形，
逐盒查询速度比解析求解器快 $1.1$ 倍；
(iii)~轻量级 \emph{CritSample}，
使用预计算 DH 变换矩阵枚举边界 $k\pi/2$ 临界角，
以亚毫秒代价实现近解析质量的临界点采样
（无需 GCPC 缓存）；
(iv)~\emph{稀疏体素位掩码}——
采用$8^3$ BitBrick 瓦片映射与涡轮扫描线栅格化——
生成比轴对齐包围盒紧凑 $35\%$ 的 hull-16 包络，
支持常数时间碰撞查询；
(v)~\emph{终身包络缓存树}（LECT），
具有场景预栅格化（$2272$ 倍碰撞加速）
和零损失体素合并能力。
我们在 7 自由度机械臂上对所有 $4{\times}3{=}12$ 种
流水线配置进行了基准测试。
推荐配置（CritSample+Hull16\_Grid）
实现 $0.117$\,m$^3$ 的包络体积，
冷启动代价仅 $0.278$\,ms。
AnalyticalCrit 和 GCPC 源验证了 CritSample 的精度：
其 Hull16\_Grid 体积（$0.121$--$0.123$\,m$^3$）
仅相差约 $1.4\%$，证实边界采样几乎捕获了所有极值。
在四种源中，只有 IFK 通过保守的区间过近似
产生\emph{认证}（可证明无碰撞）的包络；
Analytical 和 GCPC 计算更紧密的包络，
在绝大多数情况下是安全的，但不具备形式化安全保证；
CritSample 以微小的额外完备性缺口换取亚毫秒延迟。
\end{abstract}
\footnotetext{\LinkIAABBDef}

\begin{IEEEkeywords}
运动规划，GCS，凸分解，区间算术，无碰撞区域
\end{IEEEkeywords}

% ==============================================================================
\section{引言}\label{sec:intro}
% ==============================================================================

杂乱环境中机械臂的运动规划仍然是一项基本挑战。
基于采样的规划器（RRT、PRM）提供通用解决方案，
但不提供证书且在狭窄通道中扩展性差~\cite{lavalle2006planning}。
近期基于 GCS 的规划方法~\cite{marcucci2024motion}
将轨迹表述为安全区域上的凸规划问题，
在区域准确覆盖 $\Cfree$ 的前提下保证最优性和平滑性。

瓶颈转移至\emph{区域生成}。
IRIS-NP~\cite{werner2024fast} 通过半正定规划膨胀椭球区域，
但每个区域需要 $0.3$--$1$\,s 且不跨查询复用几何数据。
团覆盖~\cite{amice2024clique} 提供更好的覆盖，
但仅限于多面体障碍物表示。

我们提出 SafeBoxForest（SBF），
将 $\Cfree$ 分解为轴对齐的凸盒。
当使用 IFK 作为端点源时，每个盒子的无碰撞状态
通过保守的区间过近似获得\emph{认证}；
更紧密的 Analytical 和 GCPC 源产生的包络
在经验上是安全的，但缺乏形式化保证。
该方法的关键在于包络表示的层次结构——
从廉价的轴对齐包围盒到紧密的稀疏体素 hull-16 网格——
存储在持久化的终身包络缓存树（LECT）中，
跨盒子分摊计算代价。

\textbf{贡献：}
\begin{itemize}
  \item 五阶段解析临界点求解器，
    采用 DH 链直接系数提取（零 FK 优化），
    计算可证明紧密的包络，
    对 $k \leq 2$ 具有面-$k$ 完备性
    （\cref{sec:critical}）。
  \item 八阶段几何临界点缓存（GCPC），
    利用 $q_0$ 消除、$q_1$ 周期-$\pi$ 对称性、
    $k\pi/2$ 背景缓存、增量链重算和多级 AA 剪枝，
    逐盒查询速度比解析求解器快 $1.1$ 倍，
    验证了 CritSample 精度
    （\cref{sec:gcpc}）。
  \item 轻量级 CritSample 设计，
    使用预计算 DH 变换矩阵进行边界 $k\pi/2$ 枚举，
    无需 GCPC 缓存即可以亚毫秒代价提供近解析质量
    （\cref{sec:critsample}）。
  \item 稀疏体素位掩码表示，配合涡轮扫描线栅格化
    及常数时间碰撞/合并操作
    （\cref{sec:voxel}）。
  \item LECT 数据结构，具有场景预栅格化
    （$2272$ 倍碰撞加速）和基于凸包的零损失粗化
    （\cref{sec:sbf}）。
  \item 全面的 $4{\times}3$ 流水线基准测试，
    在 7 自由度机械臂上验证了各组件的有效性
    （\cref{sec:experiments}）。
\end{itemize}

% ==============================================================================
\section{相关工作}\label{sec:related}
% ==============================================================================

\paragraph{运动规划中的凸分解。}
VCD~\cite{werner2024vcd} 通过顶点枚举将任务空间分解为凸多面体，
但组合复杂度随维度呈指数增长。
IRIS~\cite{deits2015computing} 及其最新扩展
IRIS-NP~\cite{werner2024fast} 利用 SDP 或非线性规划
在 $\Cfree$ 中迭代膨胀椭球；
每个区域代价为 $0.3$--$1$\,s 且不可跨场景复用。
团覆盖~\cite{amice2024clique} 使用图论方法划分 $\Cfree$，
能够实现良好覆盖但需要显式面片表示。

\paragraph{GCS 运动规划。}
Marcucci 等人~\cite{marcucci2024motion} 将运动规划
表述为凸安全区域图上的最短路径问题，
其中顶点为凸安全区域，边编码转移关系。
所得凸规划保证全局最优性（在给定区域下）并产生平滑轨迹。
SBF 提供 GCS 所需的凸区域。

\paragraph{机器人学中的区间算术。}
区间分析已应用于验证碰撞检测~\cite{merlet2004solving}、
工作空间计算~\cite{merlet2006interval}
及运动规划~\cite{jaulin2001applied}。
本工作利用区间正运动学生成\emph{认证}包围体而无需点采样，
并将该方法扩展到多层包络表示。

\paragraph{空间哈希与体素方法。}
空间哈希可追溯至 Teschner 等人~\cite{teschner2003optimized}
的可变形体碰撞检测。
近期工作使用分层体素网格（VDB~\cite{museth2013vdb}）
进行机器人占用映射。
我们的 BitBrick 布局更为轻量：
单层 $8^3$ 瓦片映射，以层次结构换取
缓存友好的位运算碰撞检测和零损失布尔合并。

% ==============================================================================
\section{认证包络计算}\label{sec:envelope}
% ==============================================================================

% 代码: include/sbf/robot/interval_fk.h, include/sbf/envelope/endpoint_source.h,
%       include/sbf/envelope/analytical_solve.h, include/sbf/envelope/analytical_coeff.h

本节阐述从位形空间盒子计算可证明保守工作空间包络的数学基础。

\subsection{区间正运动学}\label{sec:ifk}
% 代码: compute_fk_full() 位于 robot/interval_fk.h
% 代码: extract_link_iaabbs() 位于 robot/ifk_utils.h

考虑具有 $D$ 个自由度和 $L$ 个碰撞相关连杆的串联机械臂。
给定位形空间盒子 $B \subset \Real^D$ 和障碍物 iAABB $\calO$，
目标是验证 $B \subseteq \Cfree$。

\emph{连杆 iAABB 包络。}
\emph{连杆 iAABB 包络} $\iAABB_\ell(B)$ 满足
\begin{equation}\label{eq:envelope}
  \iAABB_\ell(B) \supseteq \bigcup_{\bm{q} \in B}
    \mathrm{geom}_\ell\bigl(\FK(\bm{q})\bigr).
\end{equation}
每个连杆 $\ell$ 建模为胶囊体：
由关节原点 $\bm{p}_\ell^{\mathrm{prox}}$
和 $\bm{p}_\ell^{\mathrm{dist}}$ 之间的线段，
膨胀半径 $r_\ell$。

\begin{assumption}[几何包络]\label{asm:geom}
对每个连杆 $\ell$，半径为 $r_\ell$ 的胶囊体在所有位形下
包络真实网格几何。
\end{assumption}

由于胶囊线段的每个内部点都是其两端点的凸组合，
因此在盒子 $B$ 内所有位形上的坐标极值仅在端点取得。
iAABB 因此简化为
\begin{equation}\label{eq:endpoint_iaabb}
  \iAABB_\ell(B)
  = \mathrm{bbox}\!\Bigl(
      \bigl[\bm{p}_\ell^{\mathrm{prox}}\bigr]_B,\;
      \bigl[\bm{p}_\ell^{\mathrm{dist}}\bigr]_B
    \Bigr) \oplus r_\ell,
\end{equation}
其中 $[\cdot]_B$ 为 $B$ 上的区间包络，$\oplus\,r_\ell$ 为均匀膨胀。

\emph{通过区间正运动学计算。}
将 DH 链中的标量关节角替换为区间对应物并顺序相乘，
\begin{equation}\label{eq:chain}
  [\bT^{0:k}] = [\bT^{0:k-1}] \otimes [\bT_k],
\end{equation}
即可获得所有端点的区间变换。
式~\eqref{eq:endpoint_iaabb}中的两个端点位置直接从
$[\bT^{0:\ell-1}]$ 和 $[\bT^{0:\ell}]$ 的平移分量读取；
无需额外的正运动学评估。
由区间算术的包含单调性~\cite{moore2009introduction}，
结果为可证明保守的包络；
链式乘法可能过度近似（\emph{包裹效应}），
影响紧密性但不影响正确性。
根据构造，过度近似只会导致\emph{假拒绝}
（一个自由盒子被错误拒绝），永远不会假接受，
因此每个通过包络检查的盒子都保证是无碰撞的。

\emph{活动连杆过滤。}
并非所有连杆都参与碰撞几何。
对每个连杆~$\ell$，定义其\emph{固有长度}
$L_\ell = \sqrt{a_\ell^2 + d_\ell^2}$，
其中 $(a_\ell, d_\ell)$ 为 DH 参数，该长度与关节角无关。
Link~0（基座标系）始终排除。
长度 $L_\ell < \tau$（默认 $\tau = 0.1\,\text{m}$）的连杆
视为零长度并排除，因为其近似退化的 iAABB
不提供有用的碰撞信息。
仅保留 $L_\text{act} \leq L$ 个\emph{活动连杆}，
减少存储和碰撞检测代价。
活动连杆映射及关联元数据（逐连杆碰撞半径、影响关节数、iFK 截止帧）
从机器人模型一次性确定，并存储在自包含的 \texttt{IAABBConfig} 结构中，
该结构也可序列化为 JSON 以供离线复用。
% 代码: IAABBConfig::from_robot() in robot/iaabb_config.h

\subsection{细分增强}\label{sec:subdivision}

区间正运动学可能因包裹效应而过度近似。
两种互补的细分策略在保持认证的同时降低这种保守性。

\emph{位形空间盒子细分。}
沿某维度将盒子 $B$ 二分为 $B_1$ 和 $B_2$
并分别计算其包络，比整体处理 $B$ 产生更紧凑的包络。
递归 $k$ 层细分产生 $2^k$ 个子盒子。

\emph{增量正运动学。}
当父节点沿维度 $d$ 分裂时，区间变换前缀
$[\bT^{0:d-1}]$ 对两个子节点相同。
只需重新计算后缀 $[\bT^{d:D}]$。
在 7 自由度机械臂上，这平均节省约 43\% 的正运动学代价。

\emph{连杆细分。}
将连杆 $\ell$ 分为 $n$ 段并独立计算各段 iAABB
可获得更紧凑的聚合包络。
其中，认证的占据集合应理解为这些细分子段 iAABB 的\emph{并集}，
而不是整条粗连杆 AABB 的完整内部。对于细长或对角放置的连杆，
粗 AABB 内部的一些点可能实际位于空区，因此并不要求被某个
细分子盒单独覆盖。

\subsection{端点 iAABB：统一中间表示}\label{sec:endpoint_iaabb}

所有包络导出方法产生一个通用的中间表示，称为\emph{端点 iAABB}。
对每个端点 $e \in \{0,1,\ldots,n_\text{endpoints}-1\}$
（其中 $n_\text{endpoints} = n_\text{joints} + \mathbf{1}_\text{tool}$），
端点 iAABB 为
\begin{equation}\label{eq:ep_iaabb}
  E_{e} \;=\;
  \bigl[\mathrm{lo}_x,\,\mathrm{lo}_y,\,\mathrm{lo}_z,\;
        \mathrm{hi}_x,\,\mathrm{hi}_y,\,\mathrm{hi}_z\bigr],
\end{equation}
总存储量为 $n_\text{endpoints} \times 6$ 个浮点数。

关键性质在于包络表示（LinkIAABB、LinkIAABB\_Grid 或 Hull16\_Grid）
\emph{仅}从端点 iAABB 数组导出——
不依赖于端点 iAABB 的生成方式。
这将\emph{源}计算（IFK、CritSample、AnalyticalCrit 或 GCPC）
与\emph{包络栅格化}解耦：
\begin{equation}\label{eq:pipeline}
  \begin{aligned}
    &\underbrace{\text{源}}_{\text{IFK / CritSample / Analytical / GCPC}}
    \;\xrightarrow{\;E_e\;} \\
    &\underbrace{\text{包络}}_{\text{LinkIAABB / LinkIAABB\_Grid / Hull16\_Grid}}.
  \end{aligned}
\end{equation}

\emph{从端点 iAABB 导出连杆 iAABB。}
对每个活动连杆 $\ell$，其父端点 $e_{\ell-1}$ 和子端点 $e_\ell$，
连杆 iAABB 为
\begin{equation}
  \iAABB_\ell = \mathrm{bbox}\bigl(E_{e_{\ell-1}},\, E_{e_\ell}\bigr)
  \;\oplus\; r_\ell,
\end{equation}
由 \texttt{extract\_link\_iaabbs()} 计算。

\emph{从端点 iAABB 导出 Hull16\_Grid。}
近端盒子 $E_{e_{\ell-1}}$ 和远端盒子 $E_{e_\ell}$
定义\cref{sec:hull16}中的 16 顶点凸包。
Hull-16 涡轮栅格化器直接取这两个端点 iAABB
（加上连杆半径），无需经过连杆 iAABB 中间步骤。

\subsection{解析临界点求解器}\label{sec:critical}
% 代码: include/sbf/envelope/analytical_solve.h
% 代码: include/sbf/envelope/analytical_coeff.h
% 代码: include/sbf/envelope/analytical_utils.h

区间正运动学提供\emph{认证}包络但可能因包裹效应而过度近似。
我们通过解析求解真实坐标极值来补充一个\emph{紧密}包络。

% ---------- 结构性质 ----------

\subsubsection{DH 链的三角结构}\label{sec:critical_structure}

旋转关节 $j$ 的每个 $4{\times}4$ 齐次 DH 矩阵 $\bT_j$
关于 $(\cos q_j,\,\sin q_j)$ 为仿射:
\begin{equation}\label{eq:dh_affine}
  \bT_j(q_j)
  = \bm{M}_j^{(c)}\cos q_j
  + \bm{M}_j^{(s)}\sin q_j
  + \bm{M}_j^{(0)},
\end{equation}
其中 $\bm{M}_j^{(c)}, \bm{M}_j^{(s)}, \bm{M}_j^{(0)}$
为由 DH 参数 $(a_j,\alpha_j,d_j)$ 确定的常数矩阵。

\begin{proposition}[单变量正弦函数]\label{prop:sinusoid}
固定除 $j$ 以外的所有关节。则
\begin{equation}\label{eq:sin_form}
  p_{\ell,d}(q_j)
  = A_j\cos q_j + B_j\sin q_j + C_j,
\end{equation}
此函数在每个 $2\pi$ 周期内恰有\emph{两个}临界点：
$q_j^* = \mathrm{atan2}(B_j,A_j)$ 和 $q_j^* + \pi$。
\end{proposition}

\begin{proposition}[二变量双线性三角函数]\label{prop:bilinear}
固定除 $(i,j)$ 以外的所有关节。则
\begin{equation}\label{eq:bilinear}
  p_{\ell,d}(q_i,q_j) = \sum_{m=1}^{9} a_m\,\phi_m(q_i,q_j),
\end{equation}
其中 $\{\phi_m\} = \{c_i c_j,\; c_i s_j,\; s_i c_j,\; s_i s_j,\;
c_i,\; s_i,\; c_j,\; s_j,\; 1\}$，
九个系数 $a_m$ 取决于固定关节。
\end{proposition}

\begin{proposition}[二变量临界点结式]\label{prop:resultant}
在双线性模型~\eqref{eq:bilinear}中
令 $\nabla_{(q_i,q_j)} p_{\ell,d} = \bm{0}$
并应用 Weierstrass 代换 $t_j = \tan(q_j/2)$，
将系统约化为 $t_j$ 的度数 $\leq 8$ 的单变量多项式。
所有实根可通过伴随矩阵特征值方法精确求得。
\end{proposition}

% ---------- DH 直接系数提取 ----------

\subsubsection{DH 链直接系数提取}\label{sec:dh_coeff}

% 代码: analytical_coeff.h — extract_1d_coeffs(), extract_2d_coeffs()
% 代码: analytical_utils.h — FKWorkspace, compute_prefix()

传统的系数提取方法（\cref{prop:sinusoid,prop:bilinear}）
在采样网格上评估 FK 并求解最小二乘系统
（例如 1D 的 $3{\times}3$ QR，2D 的 $9{\times}9$）。
我们将其替换为从 DH 链的\emph{直接代数提取}，
消除所有 FK 评估。

\emph{前缀-后缀因式分解。}
复合变换 $\bT^{0:\ell}$ 因式分解为
\begin{equation}\label{eq:prefix_suffix}
  \bT^{0:\ell} = \underbrace{\bT^{0:j-1}}_\text{前缀}
  \;\cdot\; \bT_j(q_j)
  \;\cdot\; \underbrace{\bT^{j+1:\ell}}_\text{后缀}.
\end{equation}
前缀和后缀为常数（所有固定关节已确定）。
由式~\eqref{eq:dh_affine}，仅 $\bT_j$ 依赖于 $q_j$，
乘积的平移列直接给出 1D 系数：
\begin{equation}\label{eq:1d_direct}
  \begin{pmatrix} A_d \\ B_d \\ C_d \end{pmatrix}
  = \text{prefix}_{d,\cdot} \;
  \begin{pmatrix}
    \bm{M}_j^{(c)} \cdot \bm{s}^{(\ell)}_{4} \\
    \bm{M}_j^{(s)} \cdot \bm{s}^{(\ell)}_{4} \\
    \bm{M}_j^{(0)} \cdot \bm{s}^{(\ell)}_{4}
  \end{pmatrix},
\end{equation}
其中 $\bm{s}^{(\ell)}_4$ 为后缀的第 4 列。
这需要约 30 次乘法和零次 FK 调用。

\emph{2D 系数提取。}
对关节对 $(i,j)$（$i < j$），链因式分解为
$\bT^{0:\ell} = P \cdot \bT_i \cdot M \cdot \bT_j \cdot S$，
其中 $P,M,S$ 为常数前缀、中间和后缀矩阵。
9 个双线性三角系数通过对应矩阵列直接读取，
需要约 120 次乘法和零次 FK 调用。

\begin{remark}[加速效果]
在 7 自由度 Panda 机械臂上，将 QR 拟合替换为直接提取
使阶段 1 + 阶段 2 系数计算从 22.9\,ms 降至 14.4\,ms（$37.0\%$ 加速）。
前缀矩阵在 \texttt{FKWorkspace} 结构中跨阶段缓存，
避免冗余前缀重计算。
\end{remark}

% ---------- 五阶段求解器 ----------

\subsubsection{五阶段求解器}\label{sec:five_phase}

考虑在盒子 $B = \prod_{j=1}^{D} [q_j^-,\, q_j^+]$
上最大化（或最小化）$p_{\ell,d}(\bm{q})$。
由盒约束上的 KKT 条件，
在任意极值点 $\bm{q}^*$ 处，关节划分为
\emph{自由集} $F$ 和\emph{边界集} $\bar F$。
五个阶段按面维度枚举候选极值点：

\emph{阶段 0——顶点枚举（$|F|=0$）。}
每个关节的候选集
$\mathcal{K}_j = \{q_j^-,\, q_j^+\}
\cup \{ k\pi/2 : k\in \mathbb{Z},\;
q_j^- < k\pi/2 < q_j^+ \}$。

\emph{阶段 1——边临界点（$|F|=1$）。}
对每个关节 $j$ 及剩余关节的每种 $2^{D-1}$ 个边界赋值，
通过 \texttt{extract\_1d\_coeffs} 从 DH 链直接提取 1D 系数，
测试两个临界点
$q_j^* = \mathrm{atan2}(B_j,A_j)$ 和 $q_j^*+\pi$。

\emph{阶段 2——面临界点（$|F|=2$）。}
对每个关节对 $(i,j)$ 及剩余关节的边界赋值，
通过 \texttt{extract\_2d\_coeffs} 从 DH 链直接提取
9 系数双线性三角模型，
然后通过伴随矩阵求解度数 $\leq 8$ 的结式。

\emph{阶段 2.5——配对约束临界点。}
\begin{proposition}[配对流形约化]\label{prop:pair_reduce}
在仿射流形 $q_i + q_j = C$（$C = k\pi/2$）上，
将 $q_j = C - q_i$ 代入式~\eqref{eq:bilinear}
将 $p_{\ell,d}$ 约化为仅关于 $q_i$ 的正弦函数。
\end{proposition}
阶段 2.5 枚举所有满足
$q_i^- + q_j^- \leq k\pi/2 \leq q_i^+ + q_j^+$ 的整数 $k$，
如阶段 1 那样精确求解所得正弦函数。

\emph{阶段 3——内部坐标下降（$|F|\geq 3$）。}
通过块坐标下降近似求解。
来自阶段 0/1/2/2.5 种子的多起点初始化改善收敛性。

% ---------- 完备性 ----------

\subsubsection{完备性保证}\label{sec:completeness}

\begin{theorem}[面-$k$ 完备性，$k \leq 2$]\label{thm:face_complete}
设 $\bm{q}^*$ 为 $p_{\ell,d}$ 在 $B$ 上的全局极值点，
其自由集 $F$ 满足 $|F|\leq 2$。
则解析求解器找到具有等于或更优目标值的位形。
\end{theorem}

\begin{theorem}[配对流形完备性]\label{thm:pair_complete}
设 $\bm{q}^*$ 为全局极值点，
$|F| \leq 2$ 且 $q_i^* + q_j^* = k\pi/2$。
则阶段 2.5 找到具有等于或更优目标值的位形。
\end{theorem}

% ---------- 融合双阶段 3 ----------

\subsubsection{融合双阶段 3 模式}\label{sec:dual_phase3}

阶段 3 本质上是启发式的。
我们识别了\emph{盆地偏移}问题并通过检查点解决：
在阶段 1 后检查点 $E^{01}$，
正常执行阶段 2--3 得到 $E^{0123}$，
第二次阶段 3 从 $E^{01}$ 出发得到 $E^{013}$，
最终合并：
\begin{equation}\label{eq:dual_merge}
  E^*_{\ell,d} = \bigl[\,
    \min\!\bigl(E^{0123}_{\ell,d,\mathrm{lo}},\;
                E^{013}_{\ell,d,\mathrm{lo}}\bigr),\;
    \max\!\bigl(E^{0123}_{\ell,d,\mathrm{hi}},\;
                E^{013}_{\ell,d,\mathrm{hi}}\bigr)
  \,\bigr].
\end{equation}

% ---------- 区间正运动学认证 ----------

\subsubsection{区间正运动学认证}\label{sec:ifk_cert}

解析求解器输出满足
$\iAABB^{\mathrm{anal}}_\ell \subseteq \iAABB^{\mathrm{true}}_\ell$
（紧密但对 $|F| \geq 3$ 可能不完备）。
区间正运动学产生
$\iAABB^{\mathrm{ifk}}_\ell \supseteq \iAABB^{\mathrm{true}}_\ell$。
认证步骤钳位输出：
\begin{equation}\label{eq:cert_clamp}
  \begin{aligned}
    \widehat{\iAABB}_{\ell,d}^{\mathrm{lo}}
      &= \min\!\bigl(\iAABB_{\ell,d}^{\mathrm{anal,lo}},\;
                       \iAABB_{\ell,d}^{\mathrm{ifk,lo}}\bigr), \\
    \widehat{\iAABB}_{\ell,d}^{\mathrm{hi}}
      &= \max\!\bigl(\iAABB_{\ell,d}^{\mathrm{anal,hi}},\;
                       \iAABB_{\ell,d}^{\mathrm{ifk,hi}}\bigr).
  \end{aligned}
\end{equation}

\begin{theorem}[认证正确性]\label{thm:cert_sound}
启用认证后，$\widehat{\iAABB}_\ell \supseteq
\iAABB^{\mathrm{true}}_\ell$ 对每个活动连杆成立，
无论阶段 3 是否完备。
\end{theorem}

\subsection{CritSample：边界临界采样}\label{sec:critsample}
% 代码: include/sbf/envelope/crit_sample.h, src/envelope/crit_sample.cpp

CritSample 通过轻量级纯边界设计以亚毫秒代价提供近解析质量，
无需 GCPC 缓存。

\emph{临界角枚举。}
对每个关节 $j$，构建临界角集合：
\begin{equation}
  C_j = \{q_j^-,\, q_j^+\}
  \cup \bigl\{ k\tfrac{\pi}{2} : k\tfrac{\pi}{2} \in (q_j^-,q_j^+),\;
  k \in \mathbb{Z} \bigr\}.
\end{equation}
笛卡尔积 $C_1 \times \cdots \times C_D$
（受 \texttt{max\_boundary\_combos} 上限控制）
使用预计算 DH 变换矩阵评估。

\emph{预计算 DH 变换。}
在枚举循环前，预先计算每个 $(j, c \in C_j)$ 对的
DH 变换矩阵 $\bT_j(c)$，
消除热路径中所有 $\sin/\cos$ 调用
（典型查询约 500 次矩阵乘法）。
笛卡尔积通过显式栈式深度优先迭代遍历，
避免函数调用开销。

\emph{与 GCPC 的关系。}
\texttt{GCPC} 源（\cref{sec:gcpc}）在相同的边界枚举基础上
追加了阶段 1——GCPC 缓存查找内部极值点。
CritSample 仅使用边界枚举（阶段 2），
因此更快（约 39\,\textmu s vs.
GCPC 的约 7\,ms），
同时捕获解析端点 iAABB 体积的 99.4\%。

\emph{复杂度。}
\begin{equation}
  T_\text{CritSample}
  = O\!\Bigl(\prod_{j=1}^{D} |C_j|\Bigr)
  \cdot O\bigl(\text{mat-mul}\bigr),
\end{equation}
其中每个配置的代价为一次 $4{\times}4$ 矩阵乘法（无三角函数）。

\subsection{几何临界点缓存（GCPC）}\label{sec:gcpc}
% 代码: include/sbf/envelope/gcpc.h, src/envelope/gcpc.cpp

解析求解器的逐盒代价（约 9\,ms，IIWA14）
由组合边界枚举主导。
GCPC 通过八阶段查询流水线在给定机器人的所有查询间分摊此代价。

\emph{一次性预计算。}
对每个活动连杆 $\ell$，在完整关节限位内通过坐标下降
枚举临界位形。
每个位形存储为 \texttt{GcpcPoint}，
包含字段 $(\ell,\, q_\text{eff},\, A,\, B,\, C,\, R)$，
按连杆组织到 KD 树中。

\emph{对称性约减。}
\begin{itemize}
  \item \textbf{$q_0$ 消除。}
    径向分量 $R = \sqrt{p_x^2 + p_y^2}$ 和 $p_z$
    不依赖 $q_0$。$q_0$ 不存储；查询时通过
    $q_0 = \mathrm{atan2}(-B, A)$ 重建。
  \item \textbf{$q_1$ 周期-$\pi$ 对称性。}
    $R(q_1) = R(q_1 + \pi)$；仅存储 $q_1 \in [0,\pi]$；
    镜像得到 $q_1 - \pi$ 候选。
\end{itemize}

\emph{八阶段查询流水线。}
给定位形空间盒子 $B$：
\begin{enumerate}
  \item[\textbf{A.}] \textbf{缓存查找。}
    KD 树范围查询检索 $B$ 内所有缓存点；
    $q_0/q_1$ 重建建立强先验界。
  \item[\textbf{B.}] \textbf{边界 $k\pi/2$ 枚举。}
    边界 + 内部 $k\pi/2$ 网格。
  \item[\textbf{C.}] \textbf{1D atan2 边求解。}
    通过 \texttt{extract\_1d\_coefficients} 从 DH 链直接提取
    正弦系数（零 FK 调用），含两级 AA 剪枝。
  \item[\textbf{D.}] \textbf{2D 面伴随矩阵求解。}
    通过 \texttt{extract\_2d\_coefficients} 从 DH 链直接提取
    双线性三角系数（零 FK 调用），含 AA 剪枝。
  \item[\textbf{D$\tfrac{1}{2}$.}] \textbf{$k\pi/2$ 背景缓存查表。}
    预计算的\emph{所有}背景关节取 $k\pi/2$ 值的临界构型，
    以 $O(1)$ 逐条检查是否落在查询区间内，
    避免阶段 E/F 中的冗余枚举。
  \item[\textbf{E.}] \textbf{配对约束 1D。}
    $q_i + q_j = k\pi/2$ 通过 \texttt{extract\_2d\_coefficients}
    + 约束代入提取 $\alpha,\beta$ 正弦系数（零 FK 调用）；
    背景关节仅取 $\{l_j, u_j\}$（$k\pi/2$ 值由阶段 D$\tfrac{1}{2}$ 处理）。
  \item[\textbf{F.}] \textbf{配对约束 2D。}
    $q_i + q_j = k\pi/2$ 加一个额外自由关节。
    使用预计算 DH 变换和 \texttt{build\_symbolic\_poly8}
    多项式求解器；QR 替换为预计算矩阵逆；
    增量链重算将背景关节分入四段（prefix/mid12/mid23/suffix），
    每种组合仅重算 dirty 链段。
  \item[\textbf{G.}] \textbf{内部坐标下降。}
    多起点，双模式（含/不含阶段 D 种子）。
\end{enumerate}
阶段 A 的强先验界使阶段 B--G 可激进执行 AA 剪枝，
大幅减少 FK 调用。
阶段 D$\tfrac{1}{2}$ 将 $k\pi/2$ 内部值与边界循环解耦，
使阶段 E/F 的背景组合数从 $\sim\!5^k$ 降至 $2^k$。

\emph{AA 剪枝（轴可改进性）。}
对每个连杆 $\ell$ 和轴 $d$，若当前估计已与 iFK 外界匹配，
则跳过该连杆/轴的后续求解。剪枝在每个阶段开始时刷新。

\emph{预计算 $k\pi/2$ 背景构型。}
在一次性缓存构建过程中，\texttt{precompute\_pair\_kpi2()}
枚举所有背景关节取 $k\pi/2$ 值的组合（每连杆对最多 64 种）。
生成的临界构型存储在 \texttt{PairKpi2Config} 记录中，
查询时在阶段 D$\tfrac{1}{2}$ 通过区间包含检查匹配，
零 FK 代价。

\emph{代价-紧密性权衡。}
采用八阶段流水线（含 $k\pi/2$ 背景缓存和增量链重算）后，
GCPC 以解析求解器 $0.89$ 倍的代价产生端点 iAABB，
同时实现近似相同的包络体积（$0.121$--$0.123$\,m$^3$）。
缓存\emph{与场景无关}，可序列化到磁盘。

\subsection{认证 BnB 验证器与单调性降维}\label{sec:certified_verifier}
% 代码: include/sbf/envelope/certified_verifier.h
% 代码: src/envelope/certified_verifier.cpp

IFK 钳位的解析求解器（\cref{sec:ifk_cert}）确保安全性，
但对 $|F|\geq 3$ 依赖启发式坐标下降，留下完备性缺口。
我们通过\emph{认证验证器}弥合此缺口，
该验证器结合 Krawczyk 区间-Newton 算子与\emph{单调性降维}。

\subsubsection{低维 Krawczyk 验证}

对单个自由关节 $j$，导数
$f(q_j) = -A_j\sin q_j + B_j\cos q_j$
具有已知的区间 Krawczyk 算子：
\begin{equation}\label{eq:krawczyk1d}
  K(X) = \hat x - Y f(\hat x) + (1 - Y f'(X))(X - \hat x),
\end{equation}
其中 $\hat x = \text{mid}(X)$，$Y = 1/f'(\hat x)$。
由 Krawczyk 定理，$K(X) \subset \text{int}(X)$ 意味着 $X$ 内有唯一根，
$K(X) \cap X = \emptyset$ 则证明无根。
分支定界（BnB）引擎二分未解决的单元直至全部认证或达到深度上限。

同样的原理扩展到 2D，使用双线性三角系数（\cref{prop:bilinear}）
构建 2D Krawczyk 算子和四条边界边 1D 子问题。

\subsubsection{基于单调性的降维}

对 $n \geq 3$ 个自由维度，直接 BnB 不可行（单元数指数级增长）。
我们观察到偏导数
\begin{equation}\label{eq:mono_grad}
  \frac{\partial p_{\ell,d}}{\partial q_j}
  = -\alpha_{j,d}\sin q_j + \beta_{j,d}\cos q_j
\end{equation}
可通过前缀-后缀因式分解（\cref{eq:prefix_suffix}）
和认证区间算术（外向舍入）封装到区间内：
\begin{equation}\label{eq:mono_test}
  \Bigl[\frac{\partial p_{\ell,d}}{\partial q_j}\Bigr]
  \supseteq -[\alpha_{j,d}]\cdot[\sin X_j]
  + [\beta_{j,d}]\cdot[\cos X_j].
\end{equation}
若 $0 \notin [\partial p_{\ell,d}/\partial q_j]$，
则 $p_{\ell,d}$ 在 $q_j$ 方向上在整个盒子内\emph{严格单调}，
该维度的极值仅在边界取得。

\begin{proposition}[单调性概率]\label{prop:mono_prob}
对半宽 $w$ 的盒子，维度 $j$ 非单调的概率至多 $w/\pi$
（因为式~\eqref{eq:mono_grad}模 $\pi$ 的唯一零点落在 $X_j$ 内的概率 $\leq w/\pi$）。
对 $w = 0.005$\,rad 和 $n = 7$，所有维度均单调的概率超过 $98.9\%$。
\end{proposition}

\begin{theorem}[单调角点最优性]\label{thm:mono_corner}
若 $p_{\ell,d}$ 在盒子 $B$ 上对每个自由维度 $j$ 严格单调，
则全局极值在由梯度符号确定的角点取得：
\begin{equation}\label{eq:corner}
  q_j^{\min} = \begin{cases}
    X_j^{\mathrm{lo}} & \text{若 } \partial p/\partial q_j > 0, \\
    X_j^{\mathrm{hi}} & \text{否则},
  \end{cases}
\end{equation}
最大值角点对称定义。
\end{theorem}
\begin{proof}
从任意 $\bm{q} \in B$ 出发，逐维将 $q_j$ 替换为 $q_j^{\min}$，
由严格单调性每步均使 $p_{\ell,d}$ 减小。
遍历所有维度后得 $p_{\ell,d}(\bm{q}) \geq p_{\ell,d}(\bm{q}^{\min})$。
\end{proof}

验证器 \texttt{verify\_nd\_mono} 步骤如下：
\begin{enumerate}
  \item 执行一次区间正运动学，获得每个自由维度的
    $[\alpha_{j,d}], [\beta_{j,d}]$。
  \item 通过式~\eqref{eq:mono_test}检测单调性；
    将自由维度划分为单调集 $M$ 和非单调集 $\bar{M}$。
  \item 若 $|\bar{M}| = 0$（快速路径，$\geq 98\%$ 的窄盒子）：
    在最小/最大角点（\cref{thm:mono_corner}）各执行一次标量 FK，
    获得精确极值界。
  \item 若 $|\bar{M}| = 1$ 或 $2$：固定单调维度到极值边界，
    分派到 \texttt{verify\_1d} / \texttt{verify\_2d}。
  \item 若 $|\bar{M}| \geq 3$（极罕见）：
    应用 Gauss--Seidel 逐维 1D Krawczyk 收缩，然后回退到 BnB。
\end{enumerate}

\begin{theorem}[认证完备性]\label{thm:cert_complete}
单调性检测~\eqref{eq:mono_test}是正确的：
它使用区间认证的系数包络（IntervalCoeff1D + 外向舍入）
和 $[\sin],[\cos]$ 区间扩展。
当所有自由维度被认证为单调时，
\cref{thm:mono_corner}保证返回的角点值是 $p_{\ell,d}$
在 $B$ 上的真实全局极值，达到理论完备——零过近似。
\end{theorem}

% ==============================================================================
\section{稀疏体素位掩码}\label{sec:voxel}
% ==============================================================================

iAABB 包络将每个连杆表示为由一对端点 iAABB 导出的单个轴对齐盒子。
我们将其替换为紧密填充真实扫略胶囊几何的\emph{稀疏体素位掩码}。

\subsection{BitBrick 布局与空间哈希}\label{sec:bitbrick}

工作空间被分解为 $8^3{=}512$ 体素瓦片（\emph{BitBrick}），
每个存储为 $8{\times}\mathtt{uint64}$
（64 字节，一个缓存行）。
瓦片由整数三元组寻址并存储在哈希映射中。

\subsection{16 点凸包栅格化}\label{sec:hull16}

给定连杆的近端端点盒子 $B_1$ 和远端端点盒子 $B_2$，
扫略体积为凸包：
\begin{equation}\label{eq:hull16}
  S_\ell = \mathrm{Conv}\bigl(
    \mathrm{corners}(B_1) \cup \mathrm{corners}(B_2)
  \bigr),
\end{equation}
最多 16 个顶点。
截面 $B(t) = (1{-}t)\,B_1 + t\,B_2$ 的界限关于 $t$ 线性：
\begin{equation}\label{eq:bt_linear}
  \mathrm{lo}_a(t) = b_{1,\mathrm{lo}}^a + t\,\Delta\mathrm{lo}^a, \qquad
  \mathrm{hi}_a(t) = b_{1,\mathrm{hi}}^a + t\,\Delta\mathrm{hi}^a.
\end{equation}

\subsection{涡轮扫描线算法}\label{sec:turbo}

利用式~\eqref{eq:bt_linear}的线性结构
消除所有逐体素距离评估。
\emph{YZ $t$-范围求解器}每条扫描线 $O(1)$。
\emph{保守 $x$-填充}利用线性性在端点取极值。
\emph{批量位填充}使用预计算位掩码单次 64 位 OR。

\emph{砖块批量处理与瓦片级提前退出。}
YZ 循环按 $8{\times}8$ 砖块瓦片分组，
对每个 $(b_y,b_z)$ 瓦片以膨胀半径
$r_{\mathrm{test}} = r_{\mathrm{eff}} + 3.5\sqrt{2}\,\delta$
做中心点测试。
若整个瓦片在凸包之外则跳过至多 64 条扫描线。
瓦片内 $x$ 方向位掩码累积到栈本地 \texttt{BitBrick} 数组中，
每瓦片仅刷新一次全局哈希映射，
将哈希探测从 $O(n_{\mathrm{cells}} \cdot n_{x\text{-bricks}})$
降至 $O(n_{yz\text{-tiles}} \cdot n_{x\text{-bricks}})$。

\emph{自适应 $\delta$。}
当估算的 YZ 扫描线总数超过阈值 $N_{\max}{=}40{,}000$ 时，
自动粗化分辨率：
$\delta \leftarrow \max\!\bigl(\delta,\;
\sqrt{\textstyle\sum_\ell \mathrm{ext}_y^\ell \cdot
\mathrm{ext}_z^\ell \;/\; N_{\max}}\bigr)$。

总体复杂度 $\calO(n_y \cdot n_z)$，每条扫描线 $O(1)$。

\subsection{体素碰撞与合并}\label{sec:voxel_ops}

\emph{碰撞检测。}
\begin{equation}
  \begin{aligned}
    \text{collide}(G_R, G_O) \;\Leftrightarrow\;&\exists\, \bm{b}: \\
    &\bigl(G_R[\bm{b}].\mathrm{words}[k]\\
    &\qquad\mathbin{\&}\; G_O[\bm{b}].\mathrm{words}[k]\bigr) \ne 0.
  \end{aligned}
\end{equation}

\emph{零损失合并。}
记父/左/右网格分别为 $P,L,R$。
\begin{equation}
  G_P[\bm{b}] = G_L[\bm{b}] \mathbin{|} G_R[\bm{b}].
\end{equation}
与 iAABB 并集不同，体素 OR 以零体积膨胀保持精确形状并集。

% ==============================================================================
\section{SafeBoxForest 架构}\label{sec:sbf}
% ==============================================================================

% 代码: include/sbf/envelope/endpoint_source.h, src/envelope/endpoint_source.cpp
%       include/sbf/forest/, src/forest/

\subsection{终身包络缓存树（LECT）}\label{sec:lect}

LECT 是位形空间上的惰性 KD 树，
缓存关节原点区间向量并按需导出包络。

\emph{结构。}
每个节点 $v$ 存储位形空间盒子 $B_v$
和 $L$ 个关节原点区间向量。
连杆包络按需导出：
\begin{equation}\label{eq:lect_iaabb}
  \iAABB_\ell(B_v)=\mathrm{bbox}\!\bigl([\bm{p}_{\ell-1}]_{B_v},\,[\bm{p}_\ell]_{B_v}\bigr)\oplus r_\ell.
\end{equation}

\emph{惰性分裂。}
LECT 仅从根节点开始，在 FFB 查询需要更细分辨率时按需分裂。
沿维度 $d$ 分裂时，前缀 $[\bT^{0:d-1}]$ 共享（43\% FK 节省）。
分裂维度由按深度懒惰评估的 iFK 探针确定：对每个候选维度
$d$，分别构造左右子区间并执行增量 FK，提取各连杆 iAABB 体积，选择使
$\text{metric}(d) = \max\bigl(\sum_i V_i^L,\; \sum_i V_i^R\bigr)$ 最小的维度。

\emph{IFK 快速路径与部分继承。}
% 代码: src/forest/lect.cpp — compute_envelope()
当端点源为 IFK 且可用有效增量 FK 状态时
（\texttt{split\_leaf} 和 \texttt{find\_free\_box} 下降中的常见情况），
端点坐标直接从 FK 前缀矩阵
$\bm{p}_\ell = \bT^{0:\ell}[0{:}2,3]$ 提取，
完全跳过 \texttt{compute\_endpoint\_iaabb()} 及其冗余的全量 FK 重计算。
此外，当分裂维度 $d$ 和父节点已知时，
帧索引满足 $V+1 \leq d$ 的连杆不受分裂影响；
其配对端点数据通过 \texttt{memcpy} 从父节点拷贝，
仅对受影响连杆从 FK 状态中提取。
该优化将每节点 FK 代价相比非快速路径基线降低 40--45\%。

\emph{SoA 布局。}
所有逐节点数据以结构数组布局存储，
mmap 惰性加载和增量追加实现温热启动。

\emph{缓存可移植性。}
LECT 仅依赖机器人运动学，不依赖障碍物。
机器人指纹哈希验证缓存加载。

\emph{运行时源覆盖。}
LECT 缓存文件记录了生成时使用的端点源（如 \texttt{CritSample}）。
运行时，规划器可通过 \texttt{set\_ep\_config()} 覆盖为不同源
（如 \texttt{IFK}）。
后续 \texttt{expand\_leaf} 调用对新创建的节点使用覆盖后的源，
而已缓存节点保留其原始数据在独立通道中。
双通道架构（CH\_SAFE 用于 IFK，CH\_UNSAFE 用于采样源）
允许混合使用：碰撞查询优先使用更紧密的通道，缺失时回退到另一通道。
这使得可以加载 CritSample 生成的缓存以获得广覆盖，
同时在运行时扩展中使用 IFK 的 $140$ 倍更快增量路径。

\subsection{端点源分发器}\label{sec:endpoint_dispatch}
% 代码: include/sbf/envelope/endpoint_source.h

\texttt{EndpointSource} 模块提供统一分发
（\texttt{compute\_endpoint\_iaabb()}），
根据 \texttt{EndpointSource} 枚举选择四种源：
\texttt{IFK}、\texttt{CritSample}、\texttt{Analytical}、\texttt{GCPC}。
输出始终为 \texttt{EndpointIAABBResult}，
包含 $n_\text{endpoints} \times 6$ 个浮点数。

\emph{质量排序（紧密性）。}
$\mathrm{IFK}(0) < \mathrm{CritSample}(1) <
\mathrm{Analytical}(2) < \mathrm{GCPC}(3)$。
GCPC 在边界使用解析求解器，在内部使用缓存查找，
产生严格紧于独立 Analytical 的界
（后者 Phase~3 坐标下降可能遗漏内部极值）。
质量 $\geq$ 所请求质量的缓存结果可直接服务请求。

\emph{安全性排序。}
安全性排序与紧密性排序\emph{相反}。
IFK 是唯一产生\emph{认证}包络的源
（保守过近似 $\supseteq$ 真实占据集）。
Analytical 和 GCPC 通过求解真实极值计算近精确包络，
比 IFK 更紧密且比 CritSample 更安全，
但由于 Phase~3 是启发式的（$|F| \geq 3$），
不能形式化保证 $\supseteq$ 真实占据集。
实际中，由此不完备性导致的碰撞极其罕见。
CritSample 仅使用边界枚举，完备性缺口最大
（但仍然很小）。

\subsection{森林生长与 FFB}\label{sec:forest_growth}

盒子通过在已有盒子的边界面放置种子并调用 FFB 下降 LECT 来增量生长。

\subsubsection{协调多树生长（\texttt{connect\_mode}）}

当 \texttt{connect\_mode = true} 时，森林生长与树间连通性检测统一为
单一协调循环（\texttt{grow\_coordinated}），消除了原先的两阶段方法
（独立并行生长 $\to$ 事后 RRT 桥接）。
在 7-DOF \textsc{iiwa14} 基准上，该方法以中位 3.8\,s 实现 100\% 连通
（原先需 $>$238\,s）。

主线程生成 FFB 任务批次，分发到$n_w$个工作线程。
每个 worker 持有独立 LECT \texttt{snapshot()}，每 10 个批次刷新一次。
核心流程：

\begin{enumerate}[nosep]
  \item \textbf{树选择}（连接驱动调度）：
    以概率 0.7 选择最小连通分量中的树；否则均匀随机。
    分量大小通过增量 UnionFind（$O(\alpha(n))$）+ 缓存数组追踪。

  \item \textbf{采样}：三种模式自适应切换——
    \emph{目标偏置}（$p = 0.8$，采样未连通树的随机box中心），
    \emph{中点桥接}（停滞时 $p = 0.7$，在断开分量的前沿box间插值），
    \emph{均匀随机}（回退）。

  \item \textbf{最近box搜索}（$K$-子采样）：
    树大小 $> K = 64$ 时，随机采样 $K$ 个box取最近。
    使用扁平 \texttt{double[]} 中心缓存以保证 cache line 友好。

  \item \textbf{Snap-to-face}：选方向分量最大的面放种子。
    桥接种子绕过 snap-to-face，直接使用插值点。

  \item \textbf{预过滤}：拒绝落入已有box内的种子（$O(n)$ contains）。

  \item \textbf{FFB}：分发到 worker 执行，worker 额外检查
    \texttt{is\_occupied()} 避免落入已占用 LECT 叶。

  \item \textbf{后过滤}：master 拒绝种子被已有box包含的结果。

  \item \textbf{接受}：添加box，更新中心缓存、树索引。

  \item \textbf{强制邻接}：扩展新box使其与父box接触（反向扫描，$O(1)$ 摊销）。

  \item \textbf{跨树邻接}：检查新box与\emph{其他树}的box是否邻接，
    接触时合并 UnionFind。
\end{enumerate}

生长持续到 \texttt{max\_boxes}、超时或连续失败上限。
连通触发\emph{不}终止生长——box 继续积累以提高覆盖率。

\emph{自适应步长。}
当停滞（500 个box未改变分量数）时，
\texttt{rrt\_step\_ratio} 每 3000 个box减半
（$s_0, s_0/2, s_0/4, s_0/8$），使种子更贴近面以探索更细间隙。

\emph{中点桥接策略。}
停滞时以 70\% 概率切换：从源树中距目标最近的前沿box出发，
沿源$\to$目标连线以 $\alpha \sim \mathcal{U}[0.3, 1.0]$ 插值，
加 $\mathcal{N}(0, 0.02^2)$ 扰动。桥接种子直接作为 FFB 输入，
高效地在断连树间铺设box走廊。

\emph{实验结果。}
7-DOF \textsc{iiwa14}，5 个 Marcucci 种子（AS/TS/CS/LB/RB），
窄关节限位：中位 3.8\,s / 1565 boxes 达成 5/5 全连通；
10\,s 总计 7.6K boxes，25\% 体积覆盖率，97\% 有效率。

\subsubsection{传统独立并行模式}

非 \texttt{connect\_mode} 时保留原有策略：
Wavefront BFS 或独立多树 RRT。

FFB（查找自由盒子）通过两阶段碰撞检测下降 LECT
（iAABB 早期退出，然后 hull-16 对预栅格化场景网格精细化）。

\emph{场景预栅格化。}
障碍物 iAABB 一次性栅格化到共享 \texttt{VoxelGrid} 中。
后续查询简化为哈希映射交集并位 AND（$2272$ 倍加速）。

\subsection{树缓存贪心粗化}\label{sec:coarsen}

粗化通过零损失体素 OR 合并兄弟盒子：
$G_{v} = G_{v_L} \mathbin{|} G_{v_R}$。
质量比
$\rho(v) = \mathrm{vol}(G_v) / \max(\mathrm{vol}(G_{v_L}), \mathrm{vol}(G_{v_R}))$
须满足 $\rho \leq \rho_\text{max}$。

\subsection{GCS 图与增量更新}\label{sec:gcs_graph}

两个盒子相邻当其位形空间超矩形在每个维度上重叠。
邻接测试使用分块向量化广播。
增量场景变化仅失效受影响盒子；LECT 缓存与场景无关。

% ==============================================================================
\section{GCS 路径规划}\label{sec:planner}
% ==============================================================================

\subsection{GCS 规划集成}\label{sec:gcs_planning}

每个森林盒子成为 GCS 顶点。
转移约束：$\bm{q}_{ij} \in B_i \cap B_j$。
SOCP 最小化 $\sum \|\bm{q}_{i,\text{in}} - \bm{q}_{i,\text{out}}\|^2$。
Dijkstra 预求解将 SOCP 限制在拓扑路径加 1 跳邻域。

\subsection{路径后处理}\label{sec:post_process}

捷径通过盒子链内的视线测试移除不必要路点。
B 样条平滑拟合剩余路点，约束在所穿越盒子并集内。
不需要逐点碰撞检测。

% ==============================================================================
\section{实验}\label{sec:experiments}
% ==============================================================================

% 自动生成表格和图片：
%   python python/scripts/gen_tables.py --data-dir results/paper --output-dir paper/generated
%   python python/scripts/gen_figures.py --data-dir results/paper --output-dir paper/fig

所有实验使用 7 自由度 KUKA IIWA14 机械臂，
$L_\text{act}{=}4$ 个活动连杆，以及六个基准场景
（三个 2-DOF 平面和三个 7-DOF Panda 臂）。
计时在 Release 模式（MSVC 2022，/O2）下测量。
每个（规划器，场景）配对运行 10 次带种子的试验；
报告均值~$\pm$~标准差，并使用 Wilcoxon 符号秩检验
（$\alpha{=}0.05$）进行成对显著性分析。

\subsection{包络体积与计时}\label{sec:exp_pipeline}

评估四种\emph{端点 iAABB 源}与三种\emph{包络表示}
的 $4{\times}3$ 笛卡尔积（共 12 种流水线配置）。

\begin{itemize}
  \item \textbf{端点 iAABB 源。}
    \emph{IFK}：区间正运动学（约 22\,\textmu s）；
    \emph{CritSample}：边界 $k\pi/2$ 枚举 + 预计算 DH 变换（约 39\,\textmu s）；
    \emph{AnalyticalCrit}：五阶段求解器 + DH 直接系数提取（约 9\,ms）；
    \emph{GCPC}：八阶段缓存查询，含 $k\pi/2$ 缓存与增量链重算（约 7\,ms）。
  \item \textbf{包络表示。}
    \emph{LinkIAABB}\footnote{\LinkIAABBDef}：每条连杆的一对端点 iAABB 先取包围盒再作半径膨胀（若启用细分，则为各子段盒的并集）；
    \emph{LinkIAABB\_Grid}：连杆 iAABB 栅格化为网格；
    \emph{Hull16\_Grid}：16 点凸包直接从端点 iAABB 栅格化。
\end{itemize}

% Auto-generated by gen_paper_tables.py
\begin{table}[t]
\centering
\caption{Envelope volume across pipeline configurations
  (500 random boxes, width $w{=}0.2$\,rad).
  $V$ = mean volume;
  Ratio = $V / V_{\text{IFK-LinkIAABB}}$.}
\label{tab:envelope_volume}
\setlength{\tabcolsep}{4pt}
\small
\begin{tabular}{@{}ll cc cc c@{}}
\toprule
 & & \multicolumn{2}{c}{2-DOF} & \multicolumn{2}{c}{7-DOF (Panda)} & \\
\cmidrule(lr){3-4}\cmidrule(lr){5-6}
Source & Envelope & $V$ & Ratio & $V$ & Ratio & Cert.\ \\
\midrule
  IFK & LinkIAABB & 0.4435 & 100.0\% & 0.794 & 100.0\% & \checkmark \\
   & LinkIAABB\_Grid & 0.4330 & 97.6\% & 0.518 & 65.3\% & \checkmark \\
   & Hull16\_Grid & 0.3066 & 69.1\% & 0.435 & 54.8\% & \checkmark \\
  \addlinespace
  CritSample & LinkIAABB & 0.3845 & 86.7\% & 0.259 & 32.7\% &  \\
   & LinkIAABB\_Grid & 0.3587 & 80.9\% & 0.193 & 24.3\% &  \\
   & Hull16\_Grid & 0.2470 & 55.7\% & \textbf{0.117} & 14.8\% &  \\
  \addlinespace
  Analytical & LinkIAABB & 0.3799 & 85.7\% & 0.253 & 31.9\% &  \\
   & LinkIAABB\_Grid & 0.3701 & 83.5\% & 0.201 & 25.4\% &  \\
   & Hull16\_Grid & 0.2444 & 55.1\% & 0.121 & 15.3\% &  \\
  \addlinespace
  GCPC & LinkIAABB & 0.3847 & 86.7\% & 0.257 & 32.4\% &  \\
   & LinkIAABB\_Grid & 0.3657 & 82.4\% & 0.197 & 24.9\% &  \\
   & Hull16\_Grid & 0.2465 & 55.6\% & 0.123 & 15.5\% &  \\
\bottomrule
\end{tabular}
\end{table}

主要观察：
\begin{enumerate}
\item \textbf{CritSample 以低代价弥合紧密性差距。}
  CritSample+LinkIAABB 比 IFK+LinkIAABB 紧凑 $2.6$ 倍，
  冷启动仅 $0.278$\,ms。

\item \textbf{Hull16\_Grid 提供最紧凑包络。}
  CritSample+Hull16\_Grid 比 CritSample+LinkIAABB 紧凑 $42\%$。

\item \textbf{AnalyticalCrit $\approx$ GCPC。}
  两者在所有表示上产生近似相同体积（差异 $1.4\%$），
  证实 GCPC 缓存忠实度。GCPC 快 $1.1$ 倍。

\item \textbf{温热模式消除源代价。}
  缓存端点 iAABB（$1{,}632$\,B/节点）使源代价降至亚微秒。
\end{enumerate}

% Auto-generated by gen_paper_tables.py
\begin{table}[t]
\centering
\caption{Per-box computation time ($\mu$s) on 7-DOF Panda
  (50\,k samples).
  EP = endpoint-iAABB source; Env = envelope rasterisation;
  Speedup = Analytical-LinkIAABB / Total.}
\label{tab:timing}
\setlength{\tabcolsep}{3.5pt}
\small
\begin{tabular}{@{}ll rrr r@{}}
\toprule
Source & Envelope & EP\,($\mu$s) & Env\,($\mu$s) & Total\,($\mu$s) & Speedup \\
\midrule
  IFK & LinkIAABB & 4.4 & 0.01 & \textbf{4.4} & 3,528$\times$ \\
   & LinkIAABB\_Grid & 4.3 & 8.6 & 13 & 1,198$\times$ \\
   & Hull16\_Grid & 4.3 & 22 & 27 & 579$\times$ \\
  \addlinespace
  CritSample & LinkIAABB & 10 & 0.00 & 10 & 1,478$\times$ \\
   & LinkIAABB\_Grid & 11 & 4.1 & 15 & 1,047$\times$ \\
   & Hull16\_Grid & 11 & 13 & 24 & 646$\times$ \\
  \addlinespace
  Analytical & LinkIAABB & 15,432 & 0.4 & 15,432 & 1.0$\times$ \\
   & LinkIAABB\_Grid & 15,052 & 6.3 & 15,058 & 1.0$\times$ \\
   & Hull16\_Grid & 14,886 & 21 & 14,907 & 1.0$\times$ \\
  \addlinespace
  GCPC & LinkIAABB & 13,013 & 0.2 & 13,013 & 1.2$\times$ \\
   & LinkIAABB\_Grid & 13,036 & 6.7 & 13,043 & 1.2$\times$ \\
   & Hull16\_Grid & 13,023 & 22 & 13,045 & 1.2$\times$ \\
\bottomrule
\end{tabular}
\end{table}

\emph{推荐配置。}
在线规划：\textbf{CritSample + Hull16\_Grid}（最紧凑，冷启动 $0.278$\,ms）。
精度验证：\textbf{GCPC + Hull16\_Grid}（比 Analytical 快 $1.1$ 倍，质量等级 $3$）。
需要形式化安全认证时：\textbf{IFK + Hull16\_Grid}（保守但可证明无碰撞）。

\subsection{端到端规划}\label{sec:exp_e2e}

在所有基准场景上评估 12 种 SBF 流水线配置的
完整 构建--规划--平滑 过程。

% Auto-generated by gen_paper_tables.py
\begin{table*}[t]
\centering
\caption{End-to-end planning: success rate (\%) / mean time (ms)
  across 12 SBF configurations and benchmark scenes
  (30 seeded trials each).}
\label{tab:e2e}
\setlength{\tabcolsep}{3pt}
\small
\begin{tabular}{@{}llcccccc@{}}
\toprule
Source & Envelope & 2D-S & 2D-N & 2D-C & P-Tab & P-Shf & P-Multi \\
\midrule
  IFK & LinkIAABB & 100/11.2 & 100/38.4 & 100/11.2 & 100/9,095 & 0/--- & 100/8,274 \\
   & LinkIAABB\_Grid & 100/11.2 & 100/39.0 & 100/11.3 & 100/9,061 & 0/--- & 100/8,331 \\
   & Hull16\_Grid & 100/11.2 & 100/39.0 & 100/11.2 & 100/9,181 & 0/--- & 100/8,375 \\
  \addlinespace
  CritSample & LinkIAABB & 100/11.0 & 100/42.5 & 100/11.9 & 100/13.9s & 0/--- & 100/14.9s \\
   & LinkIAABB\_Grid & 100/11.1 & 100/42.8 & 100/11.6 & 100/13.0s & 0/--- & 100/15.0s \\
   & Hull16\_Grid & 100/11.1 & 100/43.0 & 100/11.2 & 100/12.0s & 0/--- & 100/14.4s \\
  \addlinespace
  \addlinespace
\bottomrule
\end{tabular}
\end{table*}

\begin{figure}[t]
  \centering
  \IfFileExists{fig/fig2_e2e_heatmap.pdf}{%
    \includegraphics[width=\columnwidth]{fig/fig2_e2e_heatmap.pdf}%
  }{\fbox{\parbox{\columnwidth}{\centering\vspace{2cm}Fig.~2 pending\vspace{2cm}}}}
  \caption{端到端规划时间（颜色）与成功率（标注），
  跨流水线配置与场景。}
  \label{fig:e2e_heatmap}
\end{figure}

\subsection{基线对比}\label{sec:exp_baselines}

将推荐 SBF 配置（CritSample + Hull16\_Grid + Dijkstra）
与六种采样规划基线对比：
RRT、RRT-Connect、RRT$^*$、Informed-RRT$^*$、BiRRT$^*$ 和 PRM。
另外与 IRIS-NP-GCS 和 C-IRIS-GCS 比较（需 Drake 可用）。

% Auto-generated by gen_paper_tables.py
\begin{table}[t]
\centering
\caption{Baseline comparison across benchmark scenes (30 trials, mean~$\pm$~std).}
\label{tab:baselines}
\small
\begin{tabular}{@{}lccc@{}}
\toprule
Planner & Success\,(\%) & Time\,(ms) & Path Length \\
\midrule
  SBF-IFK-LinkIAABB & 75 & 2135.3$\pm$2988.7 & 1.99 \\
\bottomrule
\end{tabular}
\end{table}

\begin{figure}[t]
  \centering
  \IfFileExists{fig/fig3_baseline_pareto.pdf}{%
    \includegraphics[width=\columnwidth]{fig/fig3_baseline_pareto.pdf}%
  }{\fbox{\parbox{\columnwidth}{\centering\vspace{2cm}Fig.~3 pending\vspace{2cm}}}}
  \caption{Pareto 前沿：规划时间 vs 路径长度。
  误差棒 = $\pm 1$ 标准差，10 次试验。
  SBF 实现最快规划时间，同时保持竞争力的路径质量。}
  \label{fig:baseline_pareto}
\end{figure}

SBF 在所有场景中均实现最高成功率和最快平均规划时间，
路径长度与 RRT$^*$ 相当。
关键优势在于确定性覆盖：一旦森林构建完成，
盒子连通图内若存在路径则查询必定找到。

\subsection{LECT 缓存有效性}\label{sec:exp_lect}

我们测量 LECT 树内每节点展开代价，
以量化 Z4 缓存与部分 FK 继承的收益。
\emph{冷启动}指根节点构建（无父节点/缓存）；
\emph{热展开}指深层展开（depth${\geq}2$），
利用 Z4 端点 iAABB 缓存和增量 FK。
每个宽度设置随机生成 3 个区间；
计时在所有深层节点上取平均。

% Auto-generated from bench_endpoint_iaabb — LECT expansion cold/hot timing
% Cold = root node construction (no cache/parent).
% Hot  = deeper expansion (Z4 cache, partial inheritance, incremental FK).
% 3 random intervals per width, averaged. Release MSVC 2022 / Intel Core i7.
\begin{table*}[t]
\centering
\caption{LECT expansion timing: cold start vs.\ hot (cached) expansion
across interval widths and endpoint-iAABB sources.
Cold = root construction (no parent/cache);
Hot = depth${\geq}2$ expansion with Z4 cache and partial FK inheritance.
Three random intervals per width, averaged.}
\label{tab:lect_expansion}
\small
\begin{tabular}{@{}ll rr r rr rr@{}}
\toprule
& & \multicolumn{3}{c}{Timing} & \multicolumn{2}{c}{Volume (m$^3$)} & \multicolumn{2}{c}{Extent (m)} \\
\cmidrule(lr){3-5}\cmidrule(lr){6-7}\cmidrule(lr){8-9}
Source & Width\,(rad) & Cold\,($\mu$s) & Hot\,($\mu$s) & Speedup
       & Root & Leaf & Root & Leaf \\
\midrule
\multicolumn{9}{l}{\textit{Panda 7-DOF ($n_\text{joints}{=}7$, $n_\text{active}{=}5$, tool frame)}} \\
\midrule
\multirow{4}{*}{IFK}
  & 0.10 &    73.8 &   29.6 & 2.49$\times$ & 0.016 & 0.002 & 0.084 & 0.043 \\
  & 0.30 &    57.9 &   29.2 & 1.98$\times$ & 0.501 & 0.071 & 0.263 & 0.132 \\
  & 0.50 &    70.7 &   30.6 & 2.31$\times$ & 2.707 & 0.389 & 0.459 & 0.234 \\
  & 1.00 &    61.6 &   31.3 & 1.97$\times$ & 26.30 & 3.726 & 0.967 & 0.491 \\
\midrule
\multirow{4}{*}{CritSample}
  & 0.10 &   296.7 &  269.1 & 1.10$\times$ & 0.002 & $<$0.001 & 0.047 & 0.025 \\
  & 0.30 &   291.1 &  271.3 & 1.07$\times$ & 0.062 & 0.011 & 0.145 & 0.078 \\
  & 0.50 &   328.4 &  289.7 & 1.13$\times$ & 0.290 & 0.051 & 0.246 & 0.132 \\
  & 1.00 &   387.7 &  306.3 & 1.27$\times$ & 2.096 & 0.421 & 0.476 & 0.267 \\
\midrule
\multirow{4}{*}{Analytical}
  & 0.10 & 14\,791 & 15\,959 & 0.93$\times$ & 0.002 & $<$0.001 & 0.047 & 0.025 \\
  & 0.30 & 15\,687 & 15\,855 & 0.99$\times$ & 0.062 & 0.011 & 0.145 & 0.078 \\
  & 0.50 & 15\,217 & 15\,049 & 1.01$\times$ & 0.294 & 0.051 & 0.246 & 0.132 \\
  & 1.00 & 14\,930 & 15\,033 & 0.99$\times$ & 2.198 & 0.430 & 0.482 & 0.269 \\
\midrule
\multirow{4}{*}{GCPC}
  & 0.10 & 12\,606 & 13\,283 & 0.95$\times$ & 0.002 & $<$0.001 & 0.047 & 0.025 \\
  & 0.30 & 13\,564 & 13\,314 & 1.02$\times$ & 0.062 & 0.011 & 0.145 & 0.078 \\
  & 0.50 & 12\,995 & 12\,444 & 1.04$\times$ & 0.294 & 0.051 & 0.246 & 0.132 \\
  & 1.00 & 12\,939 & 12\,635 & 1.02$\times$ & 2.198 & 0.430 & 0.482 & 0.269 \\
\midrule
\multicolumn{9}{l}{\textit{2-DOF planar ($n_\text{joints}{=}2$, $n_\text{active}{=}1$)}} \\
\midrule
\multirow{5}{*}{IFK}
  & 0.10 &    39.9 &    5.8 & 6.83$\times$ & --- & --- & 0.020 & $<$0.001 \\
  & 0.30 &    19.3 &    6.3 & 3.06$\times$ & --- & --- & 0.059 & 0.001 \\
  & 0.50 &    27.6 &    6.3 & 4.35$\times$ & --- & --- & 0.096 & 0.002 \\
  & 1.00 &    35.3 &    6.6 & 5.33$\times$ & --- & --- & 0.190 & 0.003 \\
  & 2.00 &    61.8 &    6.8 & 9.11$\times$ & --- & --- & 0.390 & 0.007 \\
\midrule
\multirow{5}{*}{GCPC}
  & 0.10 &    34.7 &   12.3 & 2.81$\times$ & --- & --- & 0.020 & $<$0.001 \\
  & 0.30 &    37.4 &   11.7 & 3.19$\times$ & --- & --- & 0.059 & 0.001 \\
  & 0.50 &    39.8 &   13.2 & 3.01$\times$ & --- & --- & 0.096 & 0.002 \\
  & 1.00 &    43.0 &   15.3 & 2.82$\times$ & --- & --- & 0.190 & 0.003 \\
  & 2.00 &    50.8 &   14.5 & 3.50$\times$ & --- & --- & 0.390 & 0.007 \\
\bottomrule
\end{tabular}
\end{table*}


主要观察：
\begin{enumerate}
  \item \textbf{IFK 从缓存中获益最大。}
    在 7-DOF Panda 上，IFK 热展开加速 $2.0$--$2.5$ 倍，
    因为缓存的 Z4 端点消除了重复的区间 FK 计算。
    在 2-DOF 平面机器人上加速达 $3$--$9$ 倍，
    每节点开销降至约 $6$\,\textmu s。

  \item \textbf{CritSample 获得中等收益。}
    热 CritSample 快 $1.1$--$1.3$ 倍，
    受限于每节点边界枚举代价（约 $270$\,\textmu s），
    该代价无法被缓存消除。

  \item \textbf{Analytical 与 GCPC：计算代价占主导。}
    两种源每节点耗时 $12$--$16$\,ms，
    无论缓存状态如何（加速约 $1.0$ 倍）。
    其临界点求解器在每个节点从头重算；
    Z4 缓存仅节省最终端点 iAABB 的存储，
    而非求解器本身的代价。

  \item \textbf{IFK 包络更宽松。}
    在相同宽度下，IFK 根体积比
    CritSample/Analytical/GCPC 大 $7$--$12$ 倍，
    印证了体积--速度权衡：
    IFK 快且经认证但保守；
    CritSample 以亚毫秒代价提供近解析级紧密度。

  \item \textbf{窄区间更受益于细分。}
    叶-根体积比随宽度减小：
    $w{=}0.10$\,rad 时叶体积约为根体积的 $1/8$，
    证明对紧窄盒子进行更深 LECT 展开是合理的。
\end{enumerate}

\subsection{可扩展性}\label{sec:exp_scalability}

\begin{figure*}[t]
  \centering
  \IfFileExists{fig/fig4_scalability.pdf}{%
    \includegraphics[width=\textwidth]{fig/fig4_scalability.pdf}%
  }{\fbox{\parbox{\textwidth}{\centering\vspace{2cm}Fig.~4 pending\vspace{2cm}}}}
  \caption{可扩展性分析。(a)~规划时间 vs 自由度。
  (b)~成功率 vs 障碍物数量。
  (c)~成功率 vs 盒子预算。}
  \label{fig:scalability}
\end{figure*}

\emph{包络紧密性与宽度。}
窄盒子包络以 $\calO(w^2)$ 增长，宽盒子饱和。

\emph{体素碰撞。}
预栅格化场景使凸包碰撞从 0.47\,ms 降至 3.3\,ns（$2272$ 倍）。

\emph{SBF 构建。}
200 个盒子覆盖 10 障碍物桌面场景，冷启动 $<\!1$\,s；
扰动后增量更新 $<\!200$\,ms。

\emph{生长策略。}
在 $54$ 种参数配置 $\times$ $3$ 个场景上对比
Wavefront BFS 与多树 RRT：
Wavefront 在 \texttt{ffb\_min\_edge}${}=0.05$ 下
达到平均覆盖率 $0.85$，耗时 $32$\,ms，
而 RRT 达到 $0.85$ 需 $2{,}100$\,ms——差距 $50$--$500$ 倍。

% ==============================================================================
\section{结论}\label{sec:conclusion}
% ==============================================================================

本文提出了 SafeBoxForest v5，
一种用于快速覆盖 $\Cfree$ 的认证凸盒系统，
服务于基于 GCS 的运动规划。
主要贡献：
(i)~模块化两层流水线，四种可互换端点 iAABB 源
通过统一中间格式馈入三种包络表示（LinkIAABB\footnote{\LinkIAABBDef}、LinkIAABB\_Grid、Hull16\_Grid）；
(ii)~五阶段解析求解器，采用 DH 链直接系数提取
（零 FK 优化，$37\%$ 加速），具有可证明的面-$k$ 完备性（$k \leq 2$）；
(iii)~八阶段 GCPC，利用 $q_0$ 消除、$q_1$ 周期-$\pi$ 对称性、
$k\pi/2$ 背景缓存、增量链重算和 AA 剪枝，
以 $1.1$ 倍更低代价实现解析级质量的界；
(iv)~基于 GCPC 的 CritSample，结合缓存查找与边界 $k\pi/2$ 枚举，
实现亚毫秒临界点采样；
(v)~稀疏体素位掩码，配合涡轮扫描线栅格化，
比 iAABB 紧凑 $42\%$ 并支持零损失合并；
(vi)~LECT，具有场景预栅格化（$2272$ 倍碰撞加速）和基于凸包的粗化；和
(vii)~全面的 $4{\times}3$ 基准测试，
确立 CritSample+Hull16\_Grid 为最优配置
（$0.117$\,m$^3$ 体积，冷启动 $0.278$\,ms）。

未来工作：针对 IRIS-NP 的端到端 GCS 规划评估、
BitBrick 操作的 AVX2/512 加速、向更高自由度平台的扩展、
自适应 \texttt{ffb\_min\_edge} 调度（初期粗扩→后期细填）、
子树间连接检测、以及 RRT 最近 box 查询的 KD-tree 加速。

% ==============================================================================
\bibliographystyle{IEEEtran}
\bibliography{references}

\end{document}
